% LaTeX Template for AI Problem Analysis Output
% AMC12A 2024 Problem 11 - Strategic Analysis

\documentclass[11pt]{article}
\usepackage[utf8]{inputenc}
\usepackage{amsmath, amssymb, amsthm}
\usepackage{geometry}
\usepackage{xcolor}
\usepackage[most]{tcolorbox}
\usepackage{enumitem}
\usepackage{fancyhdr}

% Page setup
\geometry{margin=1in}
\setlength{\headheight}{14.5pt}
\pagestyle{fancy}
\fancyhf{}
\rhead{AMC12/COMC Problem Analysis}
\lhead{2024 AMC 12A Problem 11}
\cfoot{\thepage}

% Color definitions
\definecolor{problemconfig}{RGB}{230, 242, 255}    % Light blue
\definecolor{strategic}{RGB}{255, 230, 242}       % Light pink
\definecolor{tactical}{RGB}{242, 255, 230}        % Light green
\definecolor{techniques}{RGB}{255, 242, 230}      % Light yellow
\definecolor{verification}{RGB}{255, 230, 230}    % Light red
\definecolor{solution}{RGB}{230, 255, 255}        % Light cyan
\definecolor{competition}{RGB}{242, 242, 255}     % Light purple
\definecolor{gold}{RGB}{218, 165, 32}

% Box styles
\newtcolorbox{problemconfigbox}{
    colback=problemconfig,
    colframe=blue,
    title=PROBLEM CONFIGURATION ANALYSIS,
    fonttitle=\bfseries,
    arc=3pt,
    boxrule=1pt,
    breakable
}

\newtcolorbox{strategicbox}{
    colback=strategic,
    colframe=purple,
    title=STRATEGIC FRAMEWORK APPLICATION,
    fonttitle=\bfseries,
    arc=3pt,
    boxrule=1pt,
    breakable
}

\newtcolorbox{tacticalbox}{
    colback=tactical,
    colframe=green,
    title=TACTICAL METHODOLOGY SELECTION,
    fonttitle=\bfseries,
    arc=3pt,
    boxrule=1pt,
    breakable
}

\newtcolorbox{techniquesbox}{
    colback=techniques,
    colframe=orange,
    title=CONVERSION TECHNIQUES APPLICATION,
    fonttitle=\bfseries,
    arc=3pt,
    boxrule=1pt,
    breakable
}

\newtcolorbox{verificationbox}{
    colback=verification,
    colframe=red,
    title=VERIFICATION STRATEGY DESIGN,
    fonttitle=\bfseries,
    arc=3pt,
    boxrule=1pt,
    breakable
}

\newtcolorbox{solutionbox}{
    colback=solution,
    colframe=cyan,
    title=SOLUTION ROADMAP,
    fonttitle=\bfseries,
    arc=3pt,
    boxrule=1pt,
    breakable
}

\newtcolorbox{competitionbox}{
    colback=competition,
    colframe=purple,
    title=COMPETITION STRATEGY INTEGRATION,
    fonttitle=\bfseries,
    arc=3pt,
    boxrule=1pt,
    breakable
}

\newtcolorbox{keyinsightbox}{
    colback=yellow!20,
    colframe=gold,
    title=KEY INSIGHT,
    fonttitle=\bfseries,
    arc=3pt,
    boxrule=2pt,
    leftrule=5pt,
    breakable
}

\newtcolorbox{warningbox}{
    colback=red!10,
    colframe=red!50,
    title=WARNING: Common Pitfall,
    fonttitle=\bfseries,
    arc=3pt,
    boxrule=2pt,
    leftrule=5pt,
    breakable
}

\newtcolorbox{tipbox}{
    colback=green!10,
    colframe=green!50,
    title=PRO TIP,
    fonttitle=\bfseries,
    arc=3pt,
    boxrule=2pt,
    leftrule=5pt,
    breakable
}

\begin{document}

\title{\textbf{AMC12A 2024 Problem 11\\Strategic Analysis}}
\author{Competition Math Framework v2.1}
\date{\today}
\maketitle

\section*{Problem Statement}

There are exactly $K$ positive integers $b$ with $5 \leq b \leq 2024$ such that the base-$b$ integer $2024_b$ is divisible by $16$ (where $16$ is in base ten). What is the sum of the digits of $K$?

\vspace{0.5em}
\textbf{Answer Choices:}\\
$\textbf{(A) }16\qquad\textbf{(B) }17\qquad\textbf{(C) }18\qquad\textbf{(D) }20\qquad\textbf{(E) }21$

\begin{problemconfigbox}
\subsection*{A. Problem Classification}
\begin{itemize}[leftmargin=*]
    \item \textbf{Problem Type}: To Find - count and compute digit sum
    \item \textbf{Competition Context}: AMC12A (75 min, 25 problems) - Problem 11 of 25
    \item \textbf{Difficulty Band}: Mid-band (P9-16) - requires number theory (bases, modular arithmetic)
    \item \textbf{Mathematical Domain}: Number Theory (base conversion, modular arithmetic, divisibility)
    \item \textbf{Estimated Time}: 3-4 minutes for students familiar with base conversion and modular arithmetic; 5-6 minutes otherwise
\end{itemize}

\subsection*{B. Problem Structure Identification}
\begin{itemize}[leftmargin=*]
    \item \textbf{Given Information}: 
    \begin{itemize}
        \item Range of bases: $5 \leq b \leq 2024$
        \item Consider the base-$b$ representation: $2024_b$
        \item Divisibility condition: $2024_b \equiv 0 \pmod{16}$ (base 10)
        \item Need to count values of $b$ satisfying this condition
        \item The digits in $2024_b$ are: 2, 0, 2, 4
    \end{itemize}
    
    \item \textbf{Constraints}: 
    \begin{itemize}
        \item $b$ must be at least 5 (since digit 4 appears, need $b > 4$)
        \item $b$ is a positive integer
        \item Upper bound: $b \leq 2024$
        \item Divisibility by 16 (a power of 2)
    \end{itemize}
    
    \item \textbf{Objective}: Find $K$ (count of valid bases), then compute the sum of digits of $K$
    
    \item \textbf{Hidden Assumptions}: 
    \begin{itemize}
        \item Base conversion formula: $2024_b = 2 \cdot b^3 + 0 \cdot b^2 + 2 \cdot b + 4 = 2b^3 + 2b + 4$
        \item Divisibility by 16 can be reduced to divisibility by 8 (factor out 2)
        \item Modular arithmetic patterns repeat with period 8
        \item Need to count using modular residues
        \item Boundary case: $b = 3$ is excluded by constraint $b \geq 5$
    \end{itemize}
    
    \item \textbf{Answer Choice Leverage}: 
    \begin{itemize}
        \item All choices are small digit sums (16-21) suggesting $K$ is in range 799-9999 or similar
        \item Specific values suggest a counting problem with clear answer
        \item The choices are close together (differ by 1-2), indicating precision matters
    \end{itemize}
\end{itemize}

\subsection*{C. Strategic Orientation}
\begin{itemize}[leftmargin=*]
    \item \textbf{Hypothesis}: Need to find all bases $b$ where $2024_b$ is divisible by 16
    \item \textbf{Conclusion}: Count such bases, then find digit sum
    \item \textbf{Notation}: 
    \begin{itemize}
        \item $2024_b$: the number 2024 in base $b$
        \item In base 10: $2024_b = 2b^3 + 2b + 4$
        \item Condition: $2b^3 + 2b + 4 \equiv 0 \pmod{16}$
        \item Simplified: $b^3 + b + 2 \equiv 0 \pmod{8}$
    \end{itemize}
    
    \item \textbf{Intuitive Approach}: 
    \begin{itemize}
        \item Convert base-$b$ to base 10 expression
        \item Set up divisibility condition as congruence
        \item Simplify modular equation
        \item Find pattern in $b \bmod 8$
        \item Count how many values in range satisfy pattern
    \end{itemize}
    
    \item \textbf{Cross-Domain Potential}: 
    \begin{itemize}
        \item \textbf{Number Theory}: Primary domain - modular arithmetic and divisibility
        \item \textbf{Algebra}: Polynomial expressions in $b$
        \item \textbf{Combinatorial Counting}: Counting residue classes in a range
        \item \textbf{Pattern Recognition}: Periodic behavior modulo 8
    \end{itemize}
\end{itemize}
\end{problemconfigbox}

\begin{strategicbox}
\subsection*{A. Psychological Strategies}
\begin{itemize}[leftmargin=*]
    \item \textbf{Mental Toughness}: This problem combines multiple concepts (bases, modular arithmetic, counting). Don't be intimidated by the multi-step nature - each step is straightforward.
    
    \item \textbf{Creativity Cultivation}: Consider testing small values of $b$ to find patterns rather than solving abstractly. Pattern recognition can guide algebraic work.
    
    \item \textbf{Iterative Refinement}: If the modular arithmetic seems complex, compute $b^3 + b + 2 \bmod 8$ for $b = 0, 1, 2, \ldots, 7$ to find which residues work.
\end{itemize}

\subsection*{B. Investigation Strategies}
\begin{itemize}[leftmargin=*]
    \item \textbf{Startup Orientation}: 
    \begin{itemize}
        \item Recognize $2024_b$ means base-$b$ notation
        \item Convert to base 10: $2024_b = 2b^3 + 0b^2 + 2b + 4$
        \item Set up divisibility: $2b^3 + 2b + 4 \equiv 0 \pmod{16}$
    \end{itemize}
    
    \item \textbf{Get Your Hands Dirty}: 
    \begin{itemize}
        \item Test $b = 5$: $2(125) + 2(5) + 4 = 250 + 10 + 4 = 264$. Is $264 \div 16 = 16.5$? No \texttimes
        \item Test $b = 6$: $2(216) + 2(6) + 4 = 432 + 12 + 4 = 448 = 16 \times 28$. Yes! \checkmark
        \item Test $b = 7$: $2(343) + 2(7) + 4 = 686 + 14 + 4 = 704 = 16 \times 44$. Yes! \checkmark
        \item Pattern emerges: check $b \bmod 8$
    \end{itemize}
    
    \item \textbf{Make It Easier}: 
    \begin{itemize}
        \item Simplify $2b^3 + 2b + 4 \equiv 0 \pmod{16}$ by dividing by 2
        \item Get $b^3 + b + 2 \equiv 0 \pmod{8}$
        \item Only need to check 8 cases: $b \equiv 0, 1, 2, \ldots, 7 \pmod{8}$
    \end{itemize}
    
    \item \textbf{Penultimate Step}: Before counting, determine which residue classes work:
    \begin{itemize}
        \item Test each $b \bmod 8$ to find: $b \equiv 3, 6, 7 \pmod{8}$
        \item These are the only solutions
        \item Count how many numbers in $[5, 2024]$ satisfy these congruences
    \end{itemize}
    
    \item \textbf{Conjecture via Small Cases}: 
    \begin{itemize}
        \item For $b \equiv 0 \pmod{8}$: $0 + 0 + 2 \equiv 2 \pmod{8}$ \texttimes
        \item For $b \equiv 1 \pmod{8}$: $1 + 1 + 2 \equiv 4 \pmod{8}$ \texttimes
        \item For $b \equiv 2 \pmod{8}$: $8 + 2 + 2 \equiv 4 \pmod{8}$ \texttimes
        \item For $b \equiv 3 \pmod{8}$: $27 + 3 + 2 = 32 \equiv 0 \pmod{8}$ \checkmark
        \item Continue for all 8 cases...
    \end{itemize}
    
    \item \textbf{Visualization}: 
    \begin{itemize}
        \item Think of numbers $0, 1, 2, \ldots, 2024$ arranged in blocks of 8
        \item Each block has exactly 3 valid values (positions 3, 6, 7)
        \item Count complete blocks: $\lfloor 2024/8 \rfloor = 253$ blocks
        \item But subtract invalid $b = 3$ (outside range $[5, 2024]$)
    \end{itemize}
\end{itemize}

\subsection*{C. Argument Construction}
\begin{itemize}[leftmargin=*]
    \item \textbf{Proof Method}: Direct computation via modular arithmetic and counting
    
    \item \textbf{Key Lemmas}: 
    \begin{itemize}
        \item \textbf{Base Conversion}: $2024_b = 2b^3 + 2b + 4$ in base 10
        \item \textbf{Divisibility Reduction}: $16 \mid 2b^3 + 2b + 4 \iff 8 \mid b^3 + b + 2$
        \item \textbf{Periodic Pattern}: $b^3 + b + 2 \bmod 8$ has period 8
        \item \textbf{Valid Residues}: $b \equiv 3, 6, 7 \pmod{8}$ are the only solutions
        \item \textbf{Counting Formula}: $K = 3 \cdot \lfloor \frac{2024}{8} \rfloor - (\text{adjustments for boundary})$
    \end{itemize}
    
    \item \textbf{Logical Structure}: 
    \begin{enumerate}
        \item Convert $2024_b$ to base 10 polynomial
        \item Set up divisibility condition modulo 16
        \item Simplify to $b^3 + b + 2 \equiv 0 \pmod{8}$
        \item Test all residue classes modulo 8
        \item Identify valid classes: $b \equiv 3, 6, 7 \pmod{8}$
        \item Count occurrences in range $[5, 2024]$
        \item Account for boundary case ($b = 3$ excluded)
        \item Compute digit sum of final count
    \end{enumerate}
\end{itemize}
\end{strategicbox}

\begin{tacticalbox}
\subsection*{A. Core Mathematical Tactics}
\begin{itemize}[leftmargin=*]
    \item \textbf{Modular Arithmetic}: Central technique - reduce polynomial modulo 8
    
    \item \textbf{Parity Arguments}: Even/odd analysis to split cases
    
    \item \textbf{Pattern Recognition}: Recognize periodicity in modular expressions
    
    \item \textbf{Systematic Case Analysis}: Check all 8 residue classes modulo 8
    
    \item \textbf{Counting Technique}: Floor division to count residue classes in a range
\end{itemize}

\subsection*{B. Domain-Specific Tactics}
\begin{itemize}[leftmargin=*]
    \item \textbf{Number Theory}: 
    \begin{itemize}
        \item Base conversion formulas
        \item Modular arithmetic reduction
        \item Divisibility rules for powers of 2
        \item Residue class counting
        \item Chinese Remainder Theorem concepts (implicit)
    \end{itemize}
    
    \item \textbf{Algebra}: 
    \begin{itemize}
        \item Polynomial evaluation
        \item Simplification of congruences
        \item Factoring (for some solution methods)
    \end{itemize}
    
    \item \textbf{Combinatorics}: 
    \begin{itemize}
        \item Counting elements in arithmetic progressions
        \item Inclusion-exclusion for boundary cases
        \item Block counting strategy
    \end{itemize}
\end{itemize}

\subsection*{C. Crossover Techniques}
\begin{itemize}[leftmargin=*]
    \item \textbf{Number Theory + Algebra}: Express divisibility as polynomial congruence
    
    \item \textbf{Number Theory + Combinatorics}: Count numbers satisfying modular conditions
    
    \item \textbf{Pattern Recognition + Systematic Testing}: Find patterns through small cases, then generalize
\end{itemize}
\end{tacticalbox}

\begin{techniquesbox}
\subsection*{A. High-Probability Techniques}
\begin{itemize}[leftmargin=*]
    \item \textbf{Primary Technique - Modular Arithmetic (H19)}: 
    \begin{itemize}
        \item \textbf{Recognition Cue}: Divisibility by 16 (power of 2), polynomial in $b$
        \item \textbf{Application}: Reduce $2b^3 + 2b + 4 \equiv 0 \pmod{16}$ to $b^3 + b + 2 \equiv 0 \pmod{8}$
        \item \textbf{Success Rate}: 90-95\% for students familiar with modular arithmetic
    \end{itemize}
    
    \item \textbf{Secondary Technique - Systematic Case Analysis}: 
    \begin{itemize}
        \item \textbf{Recognition Cue}: Small modulus (8) makes exhaustive checking feasible
        \item \textbf{Application}: Test $b = 0, 1, 2, 3, 4, 5, 6, 7$ to find pattern
        \item \textbf{Success Rate}: 85-90\% - straightforward but requires careful arithmetic
    \end{itemize}
    
    \item \textbf{Alternative - Factoring Approach}: 
    \begin{itemize}
        \item \textbf{Recognition Cue}: Can rewrite as $(b+2)(b^2-2b+5) \equiv 0 \pmod{8}$
        \item \textbf{Application}: Use factorization to analyze when product is divisible by 8
        \item \textbf{Success Rate}: 70-80\% - more elegant but requires insight
    \end{itemize}
\end{itemize}

\subsection*{B. Technique Selection Strategy}
\begin{itemize}[leftmargin=*]
    \item \textbf{Optimal Approach}: Simplify to mod 8, test all 8 cases, count valid residues
    
    \item \textbf{Backup Approach}: If modular arithmetic is unclear, compute $2b^3 + 2b + 4$ for $b = 5, 6, 7, \ldots$ until pattern emerges
    
    \item \textbf{Quick Check Approach}: Use answer choices to estimate $K$ (if digit sum is 20, then $K$ could be 758, 839, 866, etc.)
    
    \item \textbf{Time Efficiency}: 
    \begin{itemize}
        \item Modular arithmetic method: 2-3 minutes
        \item Systematic case testing: 2-3 minutes
        \item Factoring method: 3-4 minutes (if recognized)
        \item Brute force computation: 5+ minutes (not recommended)
    \end{itemize}
\end{itemize}

\subsection*{C. Technique Integration}
\begin{itemize}[leftmargin=*]
    \item \textbf{Integration Plan}:
    \begin{enumerate}
        \item Convert $2024_b$ to base 10 (30 seconds)
        \item Set up divisibility condition (20 seconds)
        \item Simplify modulo 8 (20 seconds)
        \item Test 8 cases or use parity argument (60 seconds)
        \item Identify valid residues: 3, 6, 7 (10 seconds)
        \item Count using floor division (40 seconds)
        \item Adjust for boundary ($b = 3$) (20 seconds)
        \item Compute digit sum (20 seconds)
    \end{enumerate}
    
    \item \textbf{Conflict Resolution}: If modular arithmetic is confusing, fall back to numerical testing of small values to establish pattern
\end{itemize}
\end{techniquesbox}

\begin{verificationbox}
\subsection*{A. Verification Level Selection}

Using the 7-level verification system:

\begin{itemize}[leftmargin=*]
    \item \textbf{Selected Level}: Level 4 - Standard Verification (30-60 seconds)
    \item \textbf{Rationale}: 
    \begin{itemize}
        \item Mid-band problem (P11) with multiple arithmetic steps
        \item Modular arithmetic and counting both prone to errors
        \item Can verify by checking small cases and boundary conditions
        \item 30-60 seconds appropriate for careful checking
    \end{itemize}
\end{itemize}

\subsection*{B. Verification Checklist}
\begin{itemize}[leftmargin=*]
    \item \textbf{Base Conversion Check}: 
    \begin{itemize}
        \item Verify: $2024_b = 2b^3 + 0b^2 + 2b + 4 = 2b^3 + 2b + 4$ \checkmark
        \item No terms missing
    \end{itemize}
    
    \item \textbf{Modular Arithmetic Check}:
    \begin{itemize}
        \item $2b^3 + 2b + 4 \equiv 0 \pmod{16}$
        \item Divide by 2: $b^3 + b + 2 \equiv 0 \pmod{8}$ \checkmark
        \item Valid division since $\gcd(2, 16) = 2$ and $2 \mid (2b^3 + 2b + 4)$
    \end{itemize}
    
    \item \textbf{Small Case Verification}:
    \begin{itemize}
        \item $b = 3$: $27 + 3 + 2 = 32 = 4 \times 8$ \checkmark
        \item $b = 6$: $216 + 6 + 2 = 224 = 28 \times 8$ \checkmark
        \item $b = 7$: $343 + 7 + 2 = 352 = 44 \times 8$ \checkmark
        \item $b = 11 \equiv 3$: $1331 + 11 + 2 = 1344 = 168 \times 8$ \checkmark
    \end{itemize}
    
    \item \textbf{Counting Verification}:
    \begin{itemize}
        \item $b \equiv 3 \pmod{8}$: Values are $3, 11, 19, \ldots, 2019$
        \item Count: $\frac{2019 - 3}{8} + 1 = \frac{2016}{8} + 1 = 252 + 1 = 253$
        \item But $b = 3 < 5$, so $253 - 1 = 252$ values
        \item $b \equiv 6 \pmod{8}$: Values are $6, 14, 22, \ldots, 2022$
        \item Count: $\frac{2022 - 6}{8} + 1 = \frac{2016}{8} + 1 = 253$ values
        \item $b \equiv 7 \pmod{8}$: Values are $7, 15, 23, \ldots, 2023$
        \item Count: $\frac{2023 - 7}{8} + 1 = \frac{2016}{8} + 1 = 253$ values
        \item Total: $252 + 253 + 253 = 758$
    \end{itemize}
    
    \item \textbf{Digit Sum Check}:
    \begin{itemize}
        \item $K = 758$
        \item Digit sum: $7 + 5 + 8 = 20$ \checkmark
    \end{itemize}
\end{itemize}

\subsection*{C. Confidence Calibration}
\begin{itemize}[leftmargin=*]
    \item \textbf{Target Confidence}: 90-95\%
    \item \textbf{Confidence Indicators}:
    \begin{itemize}
        \item[$\checkmark$] Modular reduction performed correctly
        \item[$\checkmark$] All 8 residue classes tested
        \item[$\checkmark$] Small cases verify pattern
        \item[$\checkmark$] Counting includes boundary adjustment
        \item[$\checkmark$] Answer matches choice (D)
    \end{itemize}
    \item \textbf{Red Flags to Watch}:
    \begin{itemize}
        \item Forgetting to exclude $b = 3$ (common error)
        \item Miscounting arithmetic progressions
        \item Off-by-one errors in counting
        \item Arithmetic errors in modular computation
    \end{itemize}
\end{itemize}
\end{verificationbox}

\begin{solutionbox}
\subsection*{A. Step-by-Step Solution}

\textbf{Method 1: Modular Arithmetic with Case Analysis (Primary)}

\textbf{Step 1: Convert base-$b$ to base 10}

The base-$b$ number $2024_b$ means:
\[
2024_b = 2 \cdot b^3 + 0 \cdot b^2 + 2 \cdot b^1 + 4 \cdot b^0 = 2b^3 + 2b + 4
\]

\textbf{Step 2: Set up divisibility condition}

We need $16 \mid 2024_b$, which means:
\[
2b^3 + 2b + 4 \equiv 0 \pmod{16}
\]

\textbf{Step 3: Simplify the congruence}

Factor out 2:
\[
2(b^3 + b + 2) \equiv 0 \pmod{16}
\]

Divide both sides by 2:
\[
b^3 + b + 2 \equiv 0 \pmod{8}
\]

\textbf{Step 4: Test all residue classes modulo 8}

We need to find which values of $b \bmod 8$ satisfy the congruence:

\begin{align*}
b \equiv 0 \pmod{8}: \quad & 0^3 + 0 + 2 \equiv 2 \pmod{8} \quad \text{\texttimes} \\
b \equiv 1 \pmod{8}: \quad & 1^3 + 1 + 2 \equiv 4 \pmod{8} \quad \text{\texttimes} \\
b \equiv 2 \pmod{8}: \quad & 2^3 + 2 + 2 \equiv 8 + 4 \equiv 4 \pmod{8} \quad \text{\texttimes} \\
b \equiv 3 \pmod{8}: \quad & 3^3 + 3 + 2 \equiv 27 + 5 \equiv 32 \equiv 0 \pmod{8} \quad \text{\checkmark} \\
b \equiv 4 \pmod{8}: \quad & 4^3 + 4 + 2 \equiv 0 + 6 \equiv 6 \pmod{8} \quad \text{\texttimes} \\
b \equiv 5 \pmod{8}: \quad & 5^3 + 5 + 2 \equiv 125 + 7 \equiv 5 + 7 \equiv 4 \pmod{8} \quad \text{\texttimes} \\
b \equiv 6 \pmod{8}: \quad & 6^3 + 6 + 2 \equiv 216 + 8 \equiv 0 + 0 \equiv 0 \pmod{8} \quad \text{\checkmark} \\
b \equiv 7 \pmod{8}: \quad & 7^3 + 7 + 2 \equiv 343 + 9 \equiv -1 + 1 \equiv 0 \pmod{8} \quad \text{\checkmark}
\end{align*}

\textbf{Step 5: Identify valid residue classes}

The valid residue classes are:
\[
b \equiv 3, 6, 7 \pmod{8}
\]

\textbf{Step 6: Count values in range $[5, 2024]$}

\textbf{For $b \equiv 3 \pmod{8}$:}
\begin{itemize}
    \item Values: $3, 11, 19, 27, \ldots, 2019$
    \item First value in range: $b = 11$ (since $b = 3 < 5$)
    \item Last value: $2019$ (since $2027 > 2024$)
    \item Count: $\frac{2019 - 11}{8} + 1 = \frac{2008}{8} + 1 = 251 + 1 = 252$
\end{itemize}

\textbf{For $b \equiv 6 \pmod{8}$:}
\begin{itemize}
    \item Values: $6, 14, 22, 30, \ldots, 2022$
    \item All values are in range $[5, 2024]$
    \item Count: $\frac{2022 - 6}{8} + 1 = \frac{2016}{8} + 1 = 252 + 1 = 253$
\end{itemize}

\textbf{For $b \equiv 7 \pmod{8}$:}
\begin{itemize}
    \item Values: $7, 15, 23, 31, \ldots, 2023$
    \item All values are in range $[5, 2024]$
    \item Count: $\frac{2023 - 7}{8} + 1 = \frac{2016}{8} + 1 = 252 + 1 = 253$
\end{itemize}

\textbf{Step 7: Compute total count}

\[
K = 252 + 253 + 253 = 758
\]

\textbf{Step 8: Find digit sum}

\[
\text{Sum of digits of } 758 = 7 + 5 + 8 = 20
\]

\textbf{Answer: (D) 20}

\subsection*{B. Alternative Approaches}

\textbf{Method 2: Block Counting}

Since the pattern repeats every 8 consecutive integers, we can think of the range $[0, 2024]$ as containing:
\[
\lfloor 2024/8 \rfloor = 253 \text{ complete blocks of 8}
\]

Each block contributes exactly 3 valid values (at positions 3, 6, 7).

Total count: $253 \times 3 = 759$

But we must exclude $b = 3$ (since $b \geq 5$):
\[
K = 759 - 1 = 758
\]

Digit sum: $7 + 5 + 8 = 20$

\textbf{Method 3: Parity Analysis}

\textbf{Case 1: $b$ is even}

Let $b = 2a$. Then:
\begin{align*}
(2a)^3 + 2a + 2 &\equiv 0 \pmod{8} \\
8a^3 + 2a + 2 &\equiv 0 \pmod{8} \\
2a + 2 &\equiv 0 \pmod{8} \\
a + 1 &\equiv 0 \pmod{4} \\
a &\equiv 3 \pmod{4}
\end{align*}

So $b = 2a = 2(4k+3) = 8k + 6$ for integer $k \geq 0$.

Thus $b \equiv 6 \pmod{8}$.

\textbf{Case 2: $b$ is odd}

Let $b = 2a + 1$. Then:
\begin{align*}
(2a+1)^3 + (2a+1) + 2 &\equiv 0 \pmod{8} \\
8a^3 + 12a^2 + 6a + 1 + 2a + 1 + 2 &\equiv 0 \pmod{8} \\
12a^2 + 8a + 4 &\equiv 0 \pmod{8} \\
4a^2 + 4 &\equiv 0 \pmod{8} \\
a^2 + 1 &\equiv 0 \pmod{2} \\
a^2 &\equiv 1 \pmod{2}
\end{align*}

This means $a$ is odd. So $a = 2m + 1$ for some integer $m \geq 0$.

Then $b = 2(2m+1) + 1 = 4m + 3$, which means $b \equiv 3 \pmod{4}$.

For $b$ odd and $b \equiv 3 \pmod{4}$:
\begin{itemize}
    \item $b = 3$: $b \equiv 3 \pmod{8}$ \checkmark
    \item $b = 7$: $b \equiv 7 \pmod{8}$ \checkmark
    \item $b = 11$: $b \equiv 3 \pmod{8}$ \checkmark
    \item $b = 15$: $b \equiv 7 \pmod{8}$ \checkmark
\end{itemize}

So odd $b$ satisfies $b \equiv 3$ or $7 \pmod{8}$.

Count as before: $K = 758$, digit sum = 20.

\textbf{Method 4: Factoring}

Notice that:
\begin{align*}
b^3 + b + 2 &\equiv 0 \pmod{8} \\
b^3 + 8 + b + 2 &\equiv 8 \pmod{8} \equiv 0 \pmod{8} \\
(b+2)(b^2-2b+4) + b + 2 &\equiv 0 \pmod{8} \\
(b+2)(b^2-2b+5) &\equiv 0 \pmod{8}
\end{align*}

Note that $b+2$ and $b^2-2b+5$ have opposite parity. For their product to be divisible by 8, one must be divisible by 8.

\textbf{If $8 \mid (b+2)$:} Then $b \equiv 6 \pmod{8}$

\textbf{If $8 \mid (b^2-2b+5)$:} Testing shows $b \equiv 3$ or $7 \pmod{8}$

Count as before: $K = 758$, digit sum = 20.

\end{solutionbox}

\begin{competitionbox}
\subsection*{A. Time Management}
\begin{itemize}[leftmargin=*]
    \item \textbf{Estimated Time}: 3-4 minutes total
    \begin{itemize}
        \item Understanding problem and base conversion: 30 seconds
        \item Setting up modular equation: 30 seconds
        \item Testing 8 cases: 60 seconds
        \item Counting values: 60 seconds
        \item Computing digit sum: 20 seconds
        \item Verification: 30 seconds
    \end{itemize}
    
    \item \textbf{Time Checkpoints}: 
    \begin{itemize}
        \item At 1 min: Should have $b^3 + b + 2 \equiv 0 \pmod{8}$
        \item At 2 min: Should have found valid residues (3, 6, 7)
        \item At 3 min: Should have counted to get $K = 758$
        \item At 4 min: Should have final answer marked
    \end{itemize}
    
    \item \textbf{Priority Level}: HIGH
    \begin{itemize}
        \item P11 is accessible with number theory background
        \item Worth full effort - systematic approach guarantees success
        \item Good point value for time invested
    \end{itemize}
\end{itemize}

\subsection*{B. Risk Assessment}
\begin{itemize}[leftmargin=*]
    \item \textbf{Confidence Level}: 90-95\%
    \begin{itemize}
        \item Systematic modular arithmetic approach
        \item Multiple verification methods available
        \item Small cases confirm pattern
        \item Answer matches choice (D)
    \end{itemize}
    
    \item \textbf{Abandonment Criteria}: 
    \begin{itemize}
        \item If after 4 minutes still making arithmetic errors $\rightarrow$ flag and move on
        \item If base conversion is confusing $\rightarrow$ test small values instead
        \item Should NOT fully abandon - pattern is discoverable
    \end{itemize}
    
    \item \textbf{Flag for Review}: MAYBE (if time permits, double-check counting)
\end{itemize}

\subsection*{C. Competition Strategy Notes}
\begin{itemize}[leftmargin=*]
    \item \textbf{Base Conversion Recognition}: The notation $2024_b$ should immediately trigger base conversion formula
    
    \item \textbf{Modulo 8 Pattern}: Divisibility by 16 for expression divisible by 2 suggests working modulo 8
    
    \item \textbf{Counting Strategy}: 
    \begin{itemize}
        \item For arithmetic progressions: use $\frac{\text{last} - \text{first}}{\text{common difference}} + 1$
        \item Always check boundary cases ($b = 3$ excluded here)
    \end{itemize}
    
    \item \textbf{Digit Sum Trick}: Computing digit sum is quick final step - be careful with arithmetic
\end{itemize}
\end{competitionbox}

\begin{keyinsightbox}
\textbf{Key Insight}: The problem reduces to finding $b$ such that $b^3 + b + 2 \equiv 0 \pmod{8}$. Since we only need to check 8 residue classes, exhaustive testing is both feasible and reliable. The pattern repeats every 8 consecutive integers with exactly 3 valid values per cycle (at positions 3, 6, 7). With 253 complete cycles in $[0, 2024]$, we get $253 \times 3 = 759$ values, minus 1 for the excluded $b = 3$, yielding $K = 758$.
\end{keyinsightbox}

\begin{warningbox}
\textbf{Warning}: The most common error is forgetting to exclude $b = 3$ from the count (since the problem requires $b \geq 5$). Also, be careful when counting arithmetic progressions - the formula is $\frac{\text{last} - \text{first}}{\text{step}} + 1$, not just the difference. Finally, when computing $b^3 \bmod 8$, don't forget to reduce intermediate results modulo 8 to keep numbers manageable (e.g., $5^3 = 125 \equiv 5 \pmod{8}$).
\end{warningbox}

\begin{tipbox}
\textbf{Pro Tip}: For AMC problems involving base conversion and divisibility, always simplify the divisibility condition as much as possible before testing cases. Here, dividing by 2 to get modulo 8 instead of modulo 16 makes the problem much easier. When counting residue classes in a range, use the ``block counting'' method: count how many complete cycles fit in the range, multiply by valid values per cycle, then adjust for boundary cases. This is faster and less error-prone than computing each arithmetic progression separately.
\end{tipbox}

\section*{Final Answer}

\boxed{\textbf{(D) } 20}

\section*{Confidence Assessment}
\begin{itemize}[leftmargin=*]
    \item \textbf{Confidence Level}: 90-95\%
    \begin{itemize}
        \item Base conversion correctly applied: $2024_b = 2b^3 + 2b + 4$
        \item Modular reduction correct: $b^3 + b + 2 \equiv 0 \pmod{8}$
        \item All 8 residue classes tested systematically
        \item Valid residues confirmed: $b \equiv 3, 6, 7 \pmod{8}$
        \item Counting verified with small cases
        \item Boundary case ($b = 3$) properly handled
        \item Multiple solution methods confirm same answer
    \end{itemize}
    
    \item \textbf{Verification Applied}: Level 4 - Standard Verification
    \begin{itemize}
        \item Checked modular arithmetic for $b = 3, 6, 7, 11$
        \item Verified counting formulas for each residue class
        \item Confirmed digit sum: $7 + 5 + 8 = 20$
        \item Block counting method yields same result
    \end{itemize}
    
    \item \textbf{Flag for Review}: Maybe (if time permits, recount to ensure no off-by-one errors)
    
    \item \textbf{Time Spent}: Approximately 3-4 minutes (optimal for P11 mid-band)
\end{itemize}

\end{document}

